% irregular-verbs.tex -- the fifty most common irregular verbs in English
%
%   make build      -> irregular-verbs.pdf (worksheet) and
%                      irregular-verbs-solutions.pdf (answer key)
%
% The fifty verbs, grouped by pattern (British English forms):
%   A-A-A  put, let, set, cut, read
%   A-B-B  have, say, get, make, think, find, tell, leave, feel, bring, keep,
%          hold, stand, hear, mean, meet, pay, sit, lead, lose, send, build,
%          understand, spend, buy
%   A-B-A  come, become, run
%   A-B-C  be, do, go, know, take, see, give, begin, write, speak, grow, fall,
%          draw, break, drive, eat, choose
%
% Note: \formrow splits its argument at commas, so table cells must not
% contain commas.
%
% Layout: a key-terms summary comes first. The parts then go from the easiest
% question type to the hardest, and each part starts on a new page:
%   1. matching (draw a line)
%   2. complete the tables
%   3. sentences with gaps
%   4. longer tasks (a story, free writing)
%
\documentclass[11pt]{article}
\usepackage{worksheet}

\begin{document}

\worksheetheader[The fifty most common irregular verbs]{Irregular Verbs}

% ===========================================================================
% Key terms
% ===========================================================================
\begin{keyterms}
  \keyterm{Infinitive}
    {the basic form of a verb, usually with \textit{to}:
     \textit{to sing}.}
  \keyterm{Past simple}
    {used for finished actions in the past:
     \textit{Yesterday I sang in the choir.}}
  \keyterm{Past participle}
    {used with \textit{have}, \textit{has} or \textit{had}:
     \textit{I have sung in the choir since 2020.}}
  \keyterm{Irregular verb}
    {a verb whose past forms do not end in \textit{-ed}.
     You have to learn them by heart.}
  \keyterm{Pattern}
    {irregular verbs fall into groups. Some have one form for everything
     (A-A-A), some have two forms the same (A-B-B), in some the past
     participle is the same as the infinitive (A-B-A) and some have three different forms (A-B-C).}

  \begin{formtable}{Infinitive, Past simple, Past participle, Pattern}
    \formrow{to hit,   hit,   hit,   A-A-A}
    \formrow{to sleep, slept, slept, A-B-B}
    \formrow{to sing,  sang,  sung,  A-B-C}
  \end{formtable}
\end{keyterms}

\clearpage

% ===========================================================================
% Part 1: matching
% ===========================================================================

% ---------------------------------------------------------------------------
\exercise{Irregular verbs: past simple}%
  {Draw a line from each infinitive to its past simple form.}

\begin{matching}[match-sep=0.9cm]{Infinitive}{Past simple}
  \matchpair{to say}{said}
  \matchpair{to mean}{meant}
  \matchpair{to stand}{stood}
  \matchpair{to understand}{understood}
  \matchpair{to do}{did}
  \matchpair{to get}{got}
  \matchpair{to make}{made}
  \matchpair{to hear}{heard}
  \matchpair{to have}{had}
  \matchpair{to find}{found}
\end{matching}

% ---------------------------------------------------------------------------
\exercise{Irregular verbs: past participle}%
  {Draw a line from each infinitive to its past participle.}

\begin{matching}[match-sep=0.9cm]{Infinitive}{Past participle}
  \matchpair{to be}{been}
  \matchpair{to go}{gone}
  \matchpair{to know}{known}
  \matchpair{to begin}{begun}
  \matchpair{to take}{taken}
  \matchpair{to give}{given}
  \matchpair{to draw}{drawn}
  \matchpair{to speak}{spoken}
  \matchpair{to write}{written}
  \matchpair{to see}{seen}
\end{matching}

\clearpage

% ===========================================================================
% Part 2: complete the tables
% ===========================================================================

% ---------------------------------------------------------------------------
\exercise{One form for everything}%
  {These verbs have the same form in all three columns. Complete the table.
   (Careful: the past forms of \textit{to read} are spelt the same but
   pronounced like \textit{red}.)}

\begin{formtable}{Infinitive, Past simple, Past participle}
  \formrow{to put,  \ans{put},  \ans{put}}
  \formrow{to let,  \ans{let},  \ans{let}}
  \formrow{to set,  \ans{set},  \ans{set}}
  \formrow{to cut,  \ans{cut},  \ans{cut}}
  \formrow{to read, \ans{read}, \ans{read}}
\end{formtable}

% ---------------------------------------------------------------------------
\exercise{Two forms the same (1)}%
  {The past simple and the past participle are the same. Complete the table.}

\begin{formtable}{Infinitive, Past simple, Past participle}
  \formrow{to have,  \ans{had},     \ans{had}}
  \formrow{to say,   \ans{said},    \ans{said}}
  \formrow{to get,   \ans{got},     \ans{got}}
  \formrow{to make,  \ans{made},    \ans{made}}
  \formrow{to think, \ans{thought}, \ans{thought}}
  \formrow{to find,  \ans{found},   \ans{found}}
  \formrow{to tell,  \ans{told},    \ans{told}}
  \formrow{to leave, \ans{left},    \ans{left}}
  \formrow{to feel,  \ans{felt},    \ans{felt}}
  \formrow{to bring, \ans{brought}, \ans{brought}}
  \formrow{to keep,  \ans{kept},    \ans{kept}}
  \formrow{to hold,  \ans{held},    \ans{held}}
  \formrow{to stand, \ans{stood},   \ans{stood}}
\end{formtable}

% ---------------------------------------------------------------------------
\exercise{Two forms the same (2)}%
  {Now some infinitives are missing, too. Complete the table.}

\begin{formtable}{Infinitive, Past simple, Past participle}
  \formrow{to hear,              \ans{heard},  \ans{heard}}
  \formrow{\ans{to mean},        meant,        \ans{meant}}
  \formrow{to meet,              \ans{met},    \ans{met}}
  \formrow{to pay,               \ans{paid},   \ans{paid}}
  \formrow{\ans{to sit},         sat,          \ans{sat}}
  \formrow{to lead,              \ans{led},    \ans{led}}
  \formrow{\ans{to lose},        lost,         \ans{lost}}
  \formrow{to send,              \ans{sent},   \ans{sent}}
  \formrow{to build,             \ans{built},  \ans{built}}
  \formrow{\ans{to understand},  understood,   \ans{understood}}
  \formrow{to spend,             \ans{spent},  \ans{spent}}
  \formrow{\ans{to buy},         bought,       \ans{bought}}
\end{formtable}

% ---------------------------------------------------------------------------
\exercise{Back to the start}%
  {The past participle is the same as the infinitive. Complete the table.}

\begin{formtable}{Infinitive, Past simple, Past participle}
  \formrow{to come,   \ans{came},   \ans{come}}
  \formrow{to become, \ans{became}, \ans{become}}
  \formrow{to run,    \ans{ran},    \ans{run}}
\end{formtable}

% ---------------------------------------------------------------------------
\exercise{Three different forms (1)}%
  {All three forms are different. Complete the table.}

\begin{formtable}{Infinitive, Past simple, Past participle}
  \formrow{to be,    \ans{was / were}, \ans{been}}
  \formrow{to do,    \ans{did},        \ans{done}}
  \formrow{to go,    \ans{went},       \ans{gone}}
  \formrow{to know,  \ans{knew},       \ans{known}}
  \formrow{to take,  \ans{took},       \ans{taken}}
  \formrow{to see,   \ans{saw},        \ans{seen}}
  \formrow{to give,  \ans{gave},       \ans{given}}
  \formrow{to begin, \ans{began},      \ans{begun}}
  \formrow{to write, \ans{wrote},      \ans{written}}
\end{formtable}

% ---------------------------------------------------------------------------
\exercise{Three different forms (2)}%
  {Complete the table. One form is given in each row.}

\begin{formtable}{Infinitive, Past simple, Past participle}
  \formrow{to speak,         \ans{spoke},  \ans{spoken}}
  \formrow{\ans{to grow},    grew,         \ans{grown}}
  \formrow{to fall,          \ans{fell},   \ans{fallen}}
  \formrow{\ans{to draw},    \ans{drew},   drawn}
  \formrow{to break,         \ans{broke},  \ans{broken}}
  \formrow{\ans{to drive},   drove,        \ans{driven}}
  \formrow{to eat,           \ans{ate},    \ans{eaten}}
  \formrow{\ans{to choose},  \ans{chose},  chosen}
\end{formtable}

\clearpage

% ===========================================================================
% Part 3: the verbs in sentences
% ===========================================================================

% ---------------------------------------------------------------------------
\exercise{Odd one out}%
  {Three verbs in each line follow the same pattern. Write the one that
   does not.}

\begin{gapexercise}
  \item put \quad cut \quad take \quad let          \hfill \gap{take}
  \item come \quad run \quad go \quad become        \hfill \gap{go}
  \item tell \quad keep \quad write \quad meet      \hfill \gap{write}
  \item speak \quad break \quad buy \quad choose    \hfill \gap{buy}
  \item know \quad grow \quad find \quad draw       \hfill \gap{find}
  \item set \quad sit \quad cut \quad read          \hfill \gap{sit}
  \item bring \quad think \quad drive \quad buy     \hfill \gap{drive}
  \item feel \quad keep \quad eat \quad mean        \hfill \gap{eat}
\end{gapexercise}

% ---------------------------------------------------------------------------
\exercise{Past simple}{Put the verb in brackets into the past simple.}

\begin{gapexercise}
  \item Last summer we \gapverb{to go}{went} to Scotland.
  \item I \gapverb{to see}{saw} your brother at the station this morning.
  \item She \gapverb{to buy}{bought} a new bike last week.
  \item We \gapverb{to eat}{ate} far too much at the party.
  \item My grandfather \gapverb{to build}{built} this house in 1965.
  \item The film \gapverb{to begin}{began} at eight o'clock.
  \item He \gapverb{to lose}{lost} his keys on the way home.
  \item I \gapverb{to think}{thought} the test was quite easy.
  \item My parents \gapverb{to meet}{met} at university.
  \item Who \gapverb{to break}{broke} the kitchen window?
  \item We \gapverb{to pay}{paid} for the tickets online.
  \item The teacher \gapverb{to tell}{told} us a funny story.
\end{gapexercise}

% ---------------------------------------------------------------------------
\exercise{Present perfect}{Put the verb in brackets into the past participle.}

\begin{gapexercise}
  \item Have you ever \gapverb{to eat}{eaten} sushi?
  \item I have \gapverb{to know}{known} her since we were children.
  \item She has \gapverb{to write}{written} three books about birds.
  \item We haven't \gapverb{to see}{seen} that film yet.
  \item Has anyone \gapverb{to take}{taken} my pen?
  \item The temperature has \gapverb{to fall}{fallen} below zero.
  \item He has \gapverb{to drive}{driven} all the way from Glasgow.
  \item They have \gapverb{to choose}{chosen} a name for the baby.
  \item I have never \gapverb{to speak}{spoken} to our new neighbours.
  \item Have you \gapverb{to do}{done} your homework?
  \item It has \gapverb{to be}{been} a very long day.
  \item My aunt has \gapverb{to give}{given} me a lot of good advice.
\end{gapexercise}

% ---------------------------------------------------------------------------
\exercise{Past simple or past participle?}%
  {Put the verb in brackets into the correct form.}

\begin{gapexercise}
  \item When I was ten, I \gapverb{to grow}{grew} ten centimetres in one year.
  \item Look how tall you have \gapverb{to grow}{grown}!
  \item She \gapverb{to leave}{left} the office at six o'clock.
  \item Has the train already \gapverb{to leave}{left}?
  \item He \gapverb{to fall}{fell} off his bike yesterday.
  \item Somebody has \gapverb{to break}{broken} my glasses.
  \item We \gapverb{to come}{came} home very late last night.
  \item I'm so glad you have \gapverb{to come}{come} to visit us.
  \item I \gapverb{to run}{ran} five kilometres this morning.
  \item Have you ever \gapverb{to run}{run} a marathon?
  \item Last year my cousin \gapverb{to become}{became} a teacher.
  \item My little sister has \gapverb{to become}{become} very good at chess.
  \item I \gapverb{to hear}{heard} a strange noise in the night.
  \item Have you \gapverb{to hear}{heard} the news?
\end{gapexercise}

% ---------------------------------------------------------------------------
\exercise{Choose the verb}%
  {Choose a verb from the box and put it into the past simple.
   Use each verb once.}

\wordbox{to bring, to feel, to hold, to keep, to lead, to send, to sit,
         to spend, to stand}

\begin{gapexercise}
  \item We \gap{sat} in the garden all afternoon.
  \item She \gap{sent} me a postcard from Rome.
  \item I \gap{felt} very tired after the long walk.
  \item The guide \gap{led} us through the old town.
  \item He \gap{held} the baby very carefully.
  \item They \gap{spent} all their money on a new car.
  \item My friend \gap{brought} a cake to the party.
  \item We \gap{stood} in the rain for an hour.
  \item She \gap{kept} all his letters in a box.
\end{gapexercise}

% ---------------------------------------------------------------------------
\exercise{Find the mistake}%
  {Each sentence has one wrong verb form. Write the correct form.}

\begin{gapexercise}
  \item Yesterday I \textit{buyed} a new jacket.            \hfill \gap{bought}
  \item She has \textit{wrote} a letter to the newspaper.    \hfill \gap{written}
  \item We \textit{goed} to the beach last weekend.          \hfill \gap{went}
  \item Have you ever \textit{drove} a lorry?                \hfill \gap{driven}
  \item He \textit{thinked} the answer was wrong.            \hfill \gap{thought}
  \item I have \textit{knew} him for years.                  \hfill \gap{known}
  \item The children have \textit{ate} all the biscuits.     \hfill \gap{eaten}
  \item In the end, he \textit{choosed} the red one.         \hfill \gap{chose}
  \item The glass has \textit{fell} off the table.           \hfill \gap{fallen}
  \item I \textit{seen} him in town yesterday.               \hfill \gap{saw}
\end{gapexercise}

\clearpage

% ===========================================================================
% Part 4: longer tasks
% ===========================================================================

% ---------------------------------------------------------------------------
\exercise{Anna's Saturday}%
  {Read the story and put the verbs in brackets into the past simple.}

{\linespread{2.0}\selectfont
Last Saturday, Anna \gapverb[2.8cm]{to get}{got} up early.
She \gapverb[2.8cm]{to have}{had} a quick breakfast and
\gapverb[2.8cm]{to put}{put} on her coat. Then she
\gapverb[2.8cm]{to drive}{drove} to the station and
\gapverb[2.8cm]{to take}{took} the train to the city.
On the train she \gapverb[2.8cm]{to read}{read} a magazine and
\gapverb[2.8cm]{to write}{wrote} a few postcards.
In the city she \gapverb[2.8cm]{to meet}{met} her old friend Sam.
They \gapverb[2.8cm]{to sit}{sat} in a café and
\gapverb[2.8cm]{to speak}{spoke} for hours.
Sam \gapverb[2.8cm]{to tell}{told} her all about his new job.
Later they \gapverb[2.8cm]{to go}{went} to a bookshop, and Anna
\gapverb[2.8cm]{to find}{found} a book that she
\gapverb[2.8cm]{to know}{knew} her brother would love, so she
\gapverb[2.8cm]{to buy}{bought} it for him.
In the evening she \gapverb[2.8cm]{to come}{came} home tired but happy
and \gapverb[2.8cm]{to fall}{fell} asleep on the sofa.
\par}

% ---------------------------------------------------------------------------
\exercise{Your turn: last weekend}%
  {Write five sentences about what you did last weekend. Use at least five
   different irregular verbs from this worksheet.}

\writinglines{6}

% ---------------------------------------------------------------------------
\exercise{Your turn: have you ever \ldots?}%
  {Write three sentences about things you have or have never done. Use
   \textit{to eat}, \textit{to see} and \textit{to be}.}

\writinglines{4}

\end{document}
